\section{Pitfalls and Recommendations}\label{sec:pitfalls}
Three simulation or design errors changed conclusions during the campaign. We describe them because each can recur in other link-adaptation studies.

\subsection{Mobility time scale}
Early versions (up to v9) moved users $v$ metres per TTI instead of $v\,T_\mathrm{tti}$. Channel dynamics were then about a thousand times too fast. Under this error the learner appeared to beat OLLA; after the correction, the v9 registry shows \NVnineWins{} wins in \NVnineCases{} cases, and the A2C policy was \NVnineVsOLLA\,\% against OLLA on goodput.

\subsection{Slot stealing}
A per-user goodput reward let a shared agent lower one user's rate to benefit others' slots. The learned policy was bimodal (\NSlotMinusThree\,\% of actions at $-3$ and \NSlotPlusTwo\,\% at $+2$), and goodput was \NSlotRL{} Mbps against \NSlotOLLA{} for OLLA. An opportunity-cost reward removed the effect (\NSlotFixRL{} against \NSlotFixOLLA{} Mbps, a $-0.14$\,\% difference that is not significant).

\subsection{PF-average timing under multi-process HARQ}
If the PF average is updated with ACKed bits at $t+\mathrm{RTT}$, a just-served user looks starved for $\mathrm{RTT}$ TTIs and is scheduled again. In a six-trace diagnostic, Baseline goodput at RTT 8 rose from \NPfOneProc{} to \NPfEightProc{} Mbps when going from one to eight processes (\NPfGain\,\%) while the Jain index fell from \NPfJainZero{} to \NPfJainEight{}. Re-running all \NPfNSchemes{} schemes with the other PF definition changed goodput by \NPfRangeLo--\NPfRangeHi\,\%, but that comparison is confounded by traces, drain and useful-bit accounting, so it is quoted only as a range. A \NEngineGap\,\% engine discontinuity had the same cause. The sign of a simple comparison depends on the definition: open-loop Baseline was \NVnineBaseGain\,\% above OLLA under delivered-throughput PF (v9) but OLLA is \NOllaBaseMthree\,\% above Baseline under scheduling-time PF in M3 (and \NOllaBaseMfour\,\% in M4). \ifitemrun{A6}{A controlled 100-trace demonstration (A6) confirms that under ACK-time averaging, going from 1 to 8 processes inflates Baseline goodput by a factor of \Asix{baseline_8p_vs_1p_acked} ($+44.7\,\%$) while degrading Jain fairness to \Asix{jain_8p_acked} ($\Delta \mathrm{Jain} = -0.097$), whereas under scheduling-time averaging the goodput ratio is \Asix{baseline_8p_vs_1p_sched} ($-2.4\,\%$) and Jain fairness is preserved at \Asix{jain_8p_sched} ($\Delta \mathrm{Jain} = +0.004$). Expectation E held across all three tested schemes, proving that the apparent multi-process HARQ gain was entirely an artefact of PF timing.}{A controlled one-trace-set demonstration (A6) is pending.}

\subsection{Withdrawn and superseded results}
Results from v7 and v8 are withdrawn. A v10 exploratory win (E4c) used overlapping traces. The v12 pre-registered replication failed Claim~B and the v13 X1 prediction failed. The v15 and v16 RTT$>$0 cases (\NRegSupCases{} cases, \NRegSupRows{} evaluations) are superseded because they used the ACK-time PF average. The registry in the supplement lists every case.

\subsection{Recommendations}
\begin{enumerate}
\item Compare against rules calibrated on the learner's cues, state the calibration budget, and give an information table such as Table~\ref{tab:info}.
\item Use common random numbers, disjoint calibration, validation and evaluation traces, and a public case registry.
\item Report seed-level intervals next to trace-level ones.
\item Report results by mobility regime, HARQ timing and observation set. The observation-set ablation in LOLLA (Table~V(a) of~\cite{wang2026lolla}, 254, 185 and 171~Mbps) shows how much a result can depend on what the learner sees.
\item State the PF-averaging definition and report the complexity of the learner.
\end{enumerate}
